\title{Επικουρικές Υπηρεσίες μη Σχετικές με τη Συχνότητα}

\section{Ευέλικτα συστήματα μεταφοράς εναλλασσόμενου ρεύματος (FACTS)}
\label{kef-02-theoria}

Πέρα από τη ρύθμιση της παραγωγής μέσω των αγορών και των εφεδρειών συχνότητας, ο Διαχειριστής
Συστήματος Μεταφοράς (ΔΣΜ) ελέγχει άμεσα τις ροές ισχύος και τις τάσεις του δικτύου μεταφοράς. Οι διατάξεις
\textbf{FACTS} είναι διατάξεις ηλεκτρονικών ισχύος που
επιτρέπουν τον ταχύ και συνεχή έλεγχο των μεγεθών που καθορίζουν τη ροή ισχύος σε δίκτυο
εναλλασσόμενου ρεύματος.

Η ενεργός ισχύς που μεταφέρεται μεταξύ δύο ζυγών μέσω γραμμής χωρίς απώλειες με επαγωγική
αντίδραση σειράς $X$ δίνεται από τη σχέση

\begin{equation}
P_{12} = \frac{V_1 V_2}{X}\,\sin(\theta_1 - \theta_2)
\label{eq:facts-powerflow}
\end{equation}

και εξαρτάται επομένως από τα μέτρα των τάσεων $V_1, V_2$, την αντίδραση σειράς $X$ και τη
διαφορά γωνιών $\theta_1 - \theta_2$. Σε
συμβατικό δίκτυο τα μεγέθη αυτά δεν ελέγχονται άμεσα, οπότε η κατανομή των ροών καθορίζεται
από τις σύνθετες αντιστάσεις των κλάδων και όχι από τις απαιτήσεις λειτουργίας. Οι διατάξεις
FACTS καθιστούν ελεγχόμενο ένα ή περισσότερα από τα μεγέθη αυτά και κατατάσσονται ανάλογα με
το μέγεθος που επηρεάζουν (Σχήμα~\ref{fig:facts}):

\begin{itemize}
\item \textbf{Παράλληλης σύνδεσης} (έλεγχος του $V$): ο στατός αντισταθμιστής αέργου ισχύος
      (Static Var Compensator, SVC) και ο στατός σύγχρονος αντισταθμιστής (Static
      Synchronous Compensator, STATCOM). Ρυθμίζουν την τάση του ζυγού μέσω έγχυσης ή
      απορρόφησης αέργου ισχύος.
\item \textbf{Σειριακής σύνδεσης} (έλεγχος του $X$): ο πυκνωτής σειράς ελεγχόμενος με
      θυρίστορ (Thyristor Controlled Series Capacitor, TCSC) και ο στατός σύγχρονος
      αντισταθμιστής σειράς (Static Synchronous Series Compensator, SSSC). Μεταβάλλουν την
      ισοδύναμη αντίδραση σειράς της γραμμής, ανακατανέμουν τις ροές ισχύος και αυξάνουν τη
      μεταφορική ικανότητα.
\item \textbf{Συνδυασμένης σύνδεσης}, σε σειρά και παράλληλα (έλεγχος και της γωνίας): ο
      ενοποιημένος ελεγκτής ροής ισχύος (Unified Power Flow Controller, UPFC) ελέγχει
      ταυτόχρονα και ανεξάρτητα την τάση και τις ροές ενεργού και αέργου ισχύος. Παρέχει τη
      μεγαλύτερη ευελιξία ελέγχου, με το υψηλότερο κόστος.
\end{itemize}

\begin{figure}[!htbp]
\centering
\includegraphics[width=\linewidth]{figures/facts.svg}
\caption{Οι τρεις κατηγορίες διατάξεων FACTS σε γραμμή δύο ζυγών. Κάθε κατηγορία επεμβαίνει
σε διαφορετικό μέγεθος της σχέσης μεταφοράς ισχύος, το οποίο σημειώνεται με έντονο χρώμα: η
παράλληλη στην τάση, η σειριακή στην αντίδραση, η συνδυασμένη και στη γωνία.}
\label{fig:facts}
\end{figure}

\subsection{Τοπολογίες διατάξεων}

Ανάλογα με την αρχή λειτουργίας τους, οι διατάξεις FACTS διακρίνονται σε διατάξεις
\emph{μεταβλητής σύνθετης αντίστασης} (εμπέδησης), με θυρίστορ, και σε διατάξεις
\emph{ελεγχόμενης πηγής τάσης}, με μετατροπείς πηγής τάσης. Τα μονογραμμικά διαγράμματα
δίνονται στο Σχήμα~\ref{fig:facts-devices}.

\begin{itemize}
\item Ο \textbf{SVC} αποτελείται από πηνίο ελεγχόμενο με θυρίστορ (Thyristor Controlled
      Reactor, TCR) και πυκνωτές μεταγόμενους με θυρίστορ (Thyristor Switched Capacitor,
      TSC), σε παράλληλη σύνδεση. Η γωνία έναυσης του TCR ρυθμίζει συνεχώς την ισοδύναμη
      επαγωγική επιδεκτικότητα, ενώ οι TSC προσθέτουν χωρητική επιδεκτικότητα σε διακριτά
      βήματα. Η διάταξη συμπεριφέρεται ως \emph{μεταβλητή επιδεκτικότητα} $B$ (susceptance).
\item Ο \textbf{STATCOM} είναι μετατροπέας πηγής τάσης (Voltage Source Converter, VSC) με
      πυκνωτή στην πλευρά συνεχούς ρεύματος (ΣΡ). Παράγει εναλλασσόμενη τάση $V_\mu$
      ελεγχόμενου μέτρου, η οποία εφαρμόζεται στο δίκτυο μέσω της αντίδρασης σύζευξης: για
      $V_\mu > V$ ο STATCOM εγχέει άεργο ισχύ, ενώ για $V_\mu < V$ απορροφά.
\item Ο \textbf{TCSC} είναι πυκνωτής σειράς με παράλληλο TCR, οπότε η ισοδύναμη αντίδραση
      σειράς της γραμμής μεταβάλλεται συνεχώς. Ο \textbf{SSSC} εγχέει τάση σε σειρά μέσω
      μετασχηματιστή σύζευξης σειράς· η εγχεόμενη τάση είναι ανεξάρτητη του ρεύματος της
      γραμμής, ενώ στον TCSC η τάση στα άκρα του πυκνωτή είναι ανάλογη του ρεύματος.
\item Ο \textbf{UPFC} αποτελείται από έναν παράλληλο και έναν σειριακό VSC με \emph{κοινό
      ζυγό ΣΡ} (DC link). Ο κοινός ζυγός επιτρέπει την ανταλλαγή \emph{ενεργού} ισχύος
      μεταξύ των δύο μετατροπέων, οπότε ο σειριακός μετατροπέας εγχέει τάση οποιουδήποτε
      μέτρου και γωνίας εντός των ονομαστικών του ορίων και ελέγχει ανεξάρτητα τις ροές
      ενεργού και αέργου ισχύος.
\end{itemize}

\begin{figure}[!htbp]
\centering
\includegraphics[width=\linewidth]{figures/facts-devices.svg}
\caption{Μονογραμμικά διαγράμματα των βασικών διατάξεων FACTS. Οι δύο πρώτες συνδέονται
παράλληλα, οι δύο επόμενες σε σειρά, ενώ ο UPFC συνδυάζει αμφότερους τους τρόπους μέσω
κοινού ζυγού ΣΡ.}
\label{fig:facts-devices}
\end{figure}

\subsection{Χαρακτηριστικές V--I των SVC και STATCOM}

Η διαφορά στην αρχή λειτουργίας αποτυπώνεται στις χαρακτηριστικές $V$--$I$ των δύο
διατάξεων παράλληλης σύνδεσης (Σχήμα~\ref{fig:facts-vi}).

Ο SVC συμπεριφέρεται ως μεταβλητή επιδεκτικότητα, οπότε $I = B\,V$ και $Q = B\,V^2$. Στο
όριο χωρητικής λειτουργίας το ρεύμα είναι ανάλογο της τάσης του ζυγού: στο 50\% της
ονομαστικής τάσης ο SVC αποδίδει το 50\% του ονομαστικού ρεύματος και το 25\% της ονομαστικής
αέργου ισχύος. Ο STATCOM, ως πηγή τάσης με έλεγχο ρεύματος, διατηρεί το ονομαστικό ρεύμα έως
πολύ χαμηλές τιμές τάσης, με περιορισμό μόνο από το μέγιστο επιτρεπόμενο ρεύμα των
ημιαγωγικών διακοπτών· η άεργος ισχύς του μειώνεται επομένως γραμμικά και όχι τετραγωνικά
με την τάση.

Κατά τη διάρκεια και αμέσως μετά από βραχυκύκλωμα, όταν η τάση είναι χαμηλή και η ανάγκη
στήριξης τάσης μέγιστη, ο STATCOM αποδίδει σημαντικά μεγαλύτερη άεργο ισχύ από SVC ίδιας
ονομαστικής ισχύος.

\begin{figure}[!htbp]
\centering
\includegraphics[width=\linewidth]{figures/facts-vi.svg}
\caption{Χαρακτηριστικά $V$--$I$ των δύο διατάξεων παράλληλης σύνδεσης. Η σκιασμένη περιοχή
είναι η περιοχή λειτουργίας και η έντονη γραμμή η χαρακτηριστική λειτουργίας με στατισμό. Τα
όρια του SVC είναι ευθείες σταθερής επιδεκτικότητας που διέρχονται από την αρχή, ενώ του
STATCOM είναι κατακόρυφες ευθείες σταθερού ρεύματος.}
\label{fig:facts-vi}
\end{figure}

Οι διατάξεις FACTS, μαζί με τα συστήματα συνεχούς ρεύματος υψηλής τάσης (High Voltage
Direct Current, HVDC), στα οποία η ροή ισχύος καθορίζεται απευθείας από τον έλεγχο των
μετατροπέων, συμπληρώνουν τα συμβατικά μέσα του συστήματος μεταφοράς: για τη ρύθμιση της τάσης
τους ρυθμιστές τάσης των γεννητριών, τα πηνία και τους πυκνωτές αντιστάθμισης και τους
μεταγωγείς λήψεων των μετασχηματιστών, και για τη διαχείριση των συμφορήσεων την ανακατανομή
παραγωγής, τις αλλαγές τοπολογίας και τους μετασχηματιστές μετατόπισης φάσης. Τα αντίστοιχα μέσα στα
δίκτυα διανομής εξετάζονται στα Κεφάλαια 5 και 6.

% TODO: συγγραφή της υπόλοιπης θεωρίας

\section{Ερωτήσεις και ασκήσεις}

\begin{enumerate}
\item \textbf{SVC και STATCOM.} Ένας SVC και ένας STATCOM έχουν ονομαστική χωρητική ισχύ
      100 Mvar. Πόση άεργο ισχύ αποδίδει ο καθένας όταν η τάση του ζυγού πέσει στα 0,7 p.u.;
      \emph{Απάντηση:} 49 Mvar και 70 Mvar αντίστοιχα.
\end{enumerate}

% TODO: παραδείγματα και λυμένες ασκήσεις (σελίδα 02b-paradeigmata.md, όπως στο Κεφάλαιο 1)
